\documentclass[12pt]{article}

\usepackage[spanish]{babel}
\addto\captionsspanish{\renewcommand{\tablename}{Tabla}}
\usepackage[utf8]{inputenc}
\usepackage[T1]{fontenc}
\usepackage{geometry}
\geometry{margin=2.5cm}
\usepackage{setspace}
\setstretch{1.15}
\usepackage{xcolor}
\definecolor{tituloazul}{RGB}{31,100,160}
\usepackage[colorlinks=true,linkcolor=tituloazul,citecolor=tituloazul,urlcolor=tituloazul]{hyperref}
\usepackage{titlesec}
\titleformat{\section}{\color{tituloazul}\large\bfseries}{\thesection}{1em}{}
\titleformat{\subsection}{\color{tituloazul}\normalsize\bfseries}{\thesubsection}{1em}{}
\titleformat{\subsubsection}{\color{tituloazul}\small\bfseries}{\thesubsubsection}{0.5em}{}
\titlespacing{\section}{0pt}{12pt}{6pt}
\titlespacing{\subsection}{0pt}{8pt}{4pt}
\usepackage{fancyhdr}
\pagestyle{fancy}
\fancyhf{}
\fancyhead[L]{\textcolor{tituloazul}{Electrónica de Potencia}}
\fancyhead[R]{\textcolor{tituloazul}{Laboratorio 3}}
\fancyfoot[C]{\thepage}
\setlength{\headheight}{15pt}
\usepackage{graphicx}
\usepackage{caption}
\usepackage{subcaption}
\usepackage{float}
\usepackage{booktabs}
\usepackage{enumitem}
\captionsetup{skip=5pt,font=small}
\setlist{topsep=5pt,itemsep=3pt,parsep=2pt}
\usepackage{amsmath}
\usepackage{siunitx}
\sisetup{per-mode=symbol,detect-all}
\usepackage[spanish,capitalize,nameinlink]{cleveref}
\crefname{figure}{figura}{figuras}
\Crefname{figure}{Figura}{Figuras}
\crefname{equation}{ecuación}{ecuaciones}
\Crefname{equation}{Ecuación}{Ecuaciones}
\crefname{table}{tabla}{tablas}
\Crefname{table}{Tabla}{Tablas}
\usepackage[siunitx]{circuitikz}
\usetikzlibrary{babel}

\newcommand{\estudiantes}{Jose Leonardo Perez Poveda, Sergio Hoyos Mejía, Alberto Mario Salazar Callejas}

\begin{document}
\vspace*{-1cm}
\begin{center}
{\Huge\textcolor{tituloazul}{\textbf{Laboratorio 3}\\Diseño y simulación de un convertidor Boost}}
\vspace{0.2cm}

{\textbf{\estudiantes}}
\vspace{0.2cm}

\textit{Ingeniería Electrónica}\\
\textit{Universidad de Antioquia}\\
\textit{Medellín-Colombia}
\end{center}
\vspace{0.3cm}

\noindent\textcolor{tituloazul}{\textbf{Resumen:}} En este laboratorio se diseñó y simuló un convertidor elevador de tensión (Boost) para elevar una entrada de \SI{10}{V} DC hasta aproximadamente \SI{17}{V} DC, con una potencia de salida cercana a \SI{15}{W} y una frecuencia de conmutación de \SI{20}{kHz}. A partir de las especificaciones se calcularon la carga, el ciclo de trabajo, el inductor y el condensador. La simulación alcanzó aproximadamente \SI{17}{V}, \SI{15}{W} y un rizado de tensión de \SI{0.82}{\percent}, inferior al límite especificado de \SI{5}{\percent}. El informe documenta además los problemas de simulación encontrados y las verificaciones pendientes antes de una implementación física.

\section{Objetivo general}
Diseñar y simular un convertidor Boost capaz de elevar una tensión de entrada de \SI{10}{V} DC hasta aproximadamente \SI{17}{V} DC, entregando una potencia cercana a \SI{15}{W} con una frecuencia de conmutación de \SI{20}{kHz}.

\section{Objetivos específicos}
\begin{itemize}
\item Calcular los parámetros principales del convertidor a partir de la tensión, potencia, frecuencia y rizado especificados.
\item Seleccionar los valores de la carga, el inductor y el condensador para la simulación.
\item Implementar en LTspice una señal PWM con el ciclo de trabajo calculado.
\item Analizar la tensión de salida, la potencia entregada y el rizado de tensión.
\item Identificar las verificaciones necesarias antes de trasladar el circuito a una implementación física.
\end{itemize}

\section{Especificaciones iniciales}
Las condiciones de diseño utilizadas se resumen en la \cref{tab:especificaciones}.
\begin{table}[H]
\centering
\caption{Especificaciones y valores utilizados en el diseño.}
\label{tab:especificaciones}
\begin{tabular}{ll}
\toprule
\textbf{Parámetro} & \textbf{Valor}\\
\midrule
Tensión de entrada, $V_{\mathrm{in}}$ & \SI{10}{V}\\
Tensión de salida, $V_{\mathrm{out}}$ & \SI{17}{V}\\
Potencia de salida, $P_{\mathrm{out}}$ & \SI{15}{W}\\
Frecuencia de conmutación, $f_s$ & \SI{20}{kHz}\\
Rizado máximo especificado & \SI{5}{\percent}\\
Carga utilizada, $R$ & \SI{19.3}{ohm}\\
Inductor, $L$ & \SI{2.7}{mH}\\
Condensador, $C$ & \SI{107}{\micro F}\\
\bottomrule
\end{tabular}
\end{table}

La resistencia de carga se calculó mediante:
\begin{equation}
R=\frac{V_{\mathrm{out}}^2}{P_{\mathrm{out}}}
\label{eq:carga}
\end{equation}
\[
R=\frac{17^2}{15}\approx\SI{19.27}{ohm}.
\]
Por lo tanto, se seleccionó el valor comercial aproximado $R=\SI{19.3}{ohm}$.

\section{Diseño del convertidor}
\subsection{Ciclo de trabajo}
Para un convertidor Boost ideal se utiliza:
\begin{equation}
V_{\mathrm{out}}=\frac{V_{\mathrm{in}}}{1-D}.
\label{eq:boost}
\end{equation}
Despejando y sustituyendo:
\begin{equation}
D=1-\frac{V_{\mathrm{in}}}{V_{\mathrm{out}}}=1-\frac{10}{17}\approx0.412.
\label{eq:duty}
\end{equation}
El valor de referencia para la señal PWM fue $D\approx\SI{41.2}{\percent}$.

\subsection{Frecuencia y señal PWM}
La frecuencia seleccionada fue $f_s=\SI{20}{kHz}$. El período correspondiente es:
\begin{equation}
T=\frac{1}{f_s}=\frac{1}{20000}=\SI{50}{\micro s}.
\label{eq:periodo}
\end{equation}
El tiempo de conducción calculado fue $T_{\mathrm{ON}}=0.412(\SI{50}{\micro s})\approx\SI{20.6}{\micro s}$. En LTspice se utilizó:
\begin{verbatim}
PULSE(0 10 0 10n 10n 20.6u 50u)
\end{verbatim}
Esta configuración genera niveles de \SI{0}{V} y \SI{10}{V}, un período de \SI{50}{\micro s} y un ciclo de trabajo aproximado de \SI{41.2}{\percent}.

\subsection{Corrientes de entrada y salida}
\begin{equation}
I_{\mathrm{out}}=\frac{P_{\mathrm{out}}}{V_{\mathrm{out}}}=\frac{15}{17}\approx\SI{0.882}{A}.
\label{eq:iout}
\end{equation}
\begin{equation}
I_{\mathrm{in}}=\frac{P_{\mathrm{out}}}{V_{\mathrm{in}}}=\frac{15}{10}=\SI{1.5}{A}.
\label{eq:iin}
\end{equation}
En un circuito real esta corriente puede ser mayor debido a las pérdidas del MOSFET, el diodo, el inductor y los demás componentes.

\subsection{Selección del inductor}
Considerando inicialmente un rizado de corriente del \SI{5}{\percent}:
\[
\Delta I_L=0.05(1.5)=\SI{0.075}{A}.
\]
Para el convertidor Boost:
\begin{equation}
L=\frac{V_{\mathrm{in}}D}{f_s\Delta I_L}=\frac{10(0.412)}{20000(0.075)}\approx\SI{2.75}{mH}.
\label{eq:inductor}
\end{equation}
Se seleccionó el valor comercial $L=\SI{2.7}{mH}$.

\subsection{Selección del condensador}
El valor teórico inicial fue cercano a \SI{21.4}{\micro F}, considerando un rizado de tensión del \SI{5}{\percent}. Durante la simulación se decidió utilizar $C=\SI{107}{\micro F}$. Este valor produce un rizado menor que el máximo permitido, aunque aumenta el tiempo necesario para alcanzar el régimen permanente.

\section{Esquema funcional}
La estructura básica descrita en el README se representa en la \cref{fig:boost}.
\begin{figure}[H]
\centering
\begin{circuitikz}[american voltages]
\draw (0,0) node[ground]{} to[battery1,l=$V_{\mathrm{in}}=\SI{10}{V}$] (0,3) to[L,l=$L=\SI{2.7}{mH}$] (3,3) -- (5,3) to[D,l=$D$] (7,3) -- (9,3) node[right]{$V_{\mathrm{out}}\approx\SI{17}{V}$};
\draw (5,3) to[short,*-] (5,1.5) to[mosfet,l_=IRFZ44N] (5,0) node[ground]{};
\draw (9,3) to[C,l=$C=\SI{107}{\micro F}$] (9,0) node[ground]{};
\draw (9,1.5) to[R,l=$R=\SI{19.3}{ohm}$] (11,1.5) node[ground]{};
\draw[->] (4,0.8) -- (5,0.8) node[midway,below]{PWM};
\end{circuitikz}
\caption{Esquema funcional del convertidor Boost simulado.}
\label{fig:boost}
\end{figure}

El MOSFET se controla mediante una señal PWM de aproximadamente \SI{20}{kHz} y \SI{41.2}{\percent} de ciclo de trabajo. El diodo, el inductor y el condensador deben seleccionarse considerando las corrientes, tensiones y pérdidas reales.

\section{Componentes utilizados en la simulación}
\begin{table}[H]
\centering
\caption{Componentes identificados en la simulación.}
\label{tab:componentes}
\begin{tabular}{lll}
\toprule
\textbf{Componente} & \textbf{Referencia o valor} & \textbf{Función}\\
\midrule
MOSFET & IRFZ44N & Interruptor de canal N\\
Diodo & 1N5819 & Diodo Schottky de conmutación\\
Inductor & \SI{2.7}{mH} & Almacenamiento de energía\\
Condensador & \SI{107}{\micro F} & Filtrado de la salida\\
Carga & \SI{19.3}{ohm} & Demanda de potencia\\
\bottomrule
\end{tabular}
\end{table}

El 1N5819 se utilizó en la primera simulación por su baja caída de tensión y sus características de conmutación. Para la implementación física se debe verificar que su tensión inversa, corriente y disipación sean suficientes.

\section{Resultados de simulación y análisis}
\subsection{Tensión de salida y rizado}
En la simulación se observaron aproximadamente $V_{\max}=\SI{17.15}{V}$ y $V_{\min}=\SI{17.01}{V}$. Por tanto:
\begin{equation}
\Delta V=17.15-17.01=\SI{0.14}{V}.
\label{eq:rizado}
\end{equation}
\begin{equation}
\%\,\mathrm{Ripple}=\frac{0.14}{17}(100)\approx\SI{0.82}{\percent}.
\label{eq:porcentaje_rizado}
\end{equation}
El resultado es inferior al límite especificado de \SI{5}{\percent}. El condensador proporciona una salida estable, a costa de un mayor tiempo de establecimiento.

\subsection{Potencia de salida}
\begin{equation}
P_{\mathrm{out}}=V_{\mathrm{out}}I_{\mathrm{out}}=\frac{V_{\mathrm{out}}^2}{R}.
\label{eq:pout}
\end{equation}
Para aproximadamente \SI{17}{V} y \SI{19.3}{ohm}:
\[
P_{\mathrm{out}}=\frac{17^2}{19.3}\approx\SI{14.97}{W}.
\]
Este valor es cercano a los \SI{15}{W} especificados.

\subsection{Factor de potencia y eficiencia}
El factor de potencia no es equivalente a la eficiencia. En este proyecto la entrada es de \SI{10}{VDC}, por lo que el concepto de factor de potencia no se aplica de la misma forma que en un sistema alimentado directamente desde una red de corriente alterna. La eficiencia debe calcularse mediante:
\begin{equation}
\eta=\frac{P_{\mathrm{out}}}{P_{\mathrm{in}}}\times100\%.
\label{eq:eficiencia}
\end{equation}
El README no proporciona la potencia de entrada simulada; por tanto, no es posible reportar un valor numérico de eficiencia sin inventar datos. También quedan pendientes la medición del rizado de corriente y las pérdidas de los componentes.

\section{Problemas encontrados y correcciones}
\begin{table}[H]
\centering
\caption{Problemas identificados durante la simulación.}
\label{tab:problemas}
\begin{tabular}{p{4.2cm}p{8.5cm}}
\toprule
\textbf{Problema} & \textbf{Corrección aplicada}\\
\midrule
MOSFET de canal incorrecto & Se reemplazó el MOSFET de canal P por el MOSFET de canal N IRFZ44N.\\
PWM incorrecto & Se cambió la fuente inicial por \texttt{PULSE(0 10 0 10n 10n 20.6u 50u)}, correspondiente al ciclo calculado.\\
Tiempo de simulación insuficiente & Se aumentó el tiempo de simulación; se propuso \texttt{.tran 30m} para observar el arranque y el régimen permanente.\\
Tensión de salida inesperada & Los valores cercanos a \SI{0.3}{V} y \SI{11.2}{V} se asociaron a la configuración del MOSFET y posteriormente a la señal de control.\\
\bottomrule
\end{tabular}
\end{table}

\section{Consideraciones para la implementación física}
Antes de llevar el circuito a una board se deben realizar las verificaciones de la \cref{tab:verificaciones}. Estas tareas no se reportan como resultados experimentales porque el README las identifica como trabajo pendiente.
\begin{table}[H]
\centering
\caption{Verificaciones pendientes antes de la implementación física.}
\label{tab:verificaciones}
\begin{tabular}{p{8.5cm}p{4.2cm}}
\toprule
\textbf{Verificación} & \textbf{Estado}\\
\midrule
Confirmar tensión y corriente reales de entrada & Pendiente\\
Medir potencia de entrada y calcular eficiencia & Pendiente\\
Verificar rizado de corriente del inductor & Pendiente\\
Verificar corriente máxima del MOSFET & Pendiente\\
Verificar corriente y tensión inversa del diodo & Pendiente\\
Confirmar corriente de saturación del inductor & Pendiente\\
Verificar potencia disipada y temperatura del MOSFET & Pendiente\\
Seleccionar los componentes físicos definitivos & Pendiente\\
Implementar el circuito y medir sus variables principales & Pendiente\\
\bottomrule
\end{tabular}
\end{table}

El inductor debe soportar una corriente superior a la corriente máxima del circuito, considerando su corriente de saturación, resistencia serie y pérdidas. El condensador debe tener una tensión nominal superior a \SI{17}{V}; el README recomienda como referencia al menos \SI{25}{V}, sujeto al diseño final, y verificar también su corriente de rizado.

Para el IRFZ44N se deben comprobar $V_{\mathrm{DS}}$, corriente máxima, $R_{\mathrm{DS(on)}}$, disipación térmica y $V_{\mathrm{GS}}$. También se recomienda utilizar una resistencia entre el PWM y la compuerta y una resistencia de \emph{pull-down} entre compuerta y fuente. El cableado de potencia debe ser corto y mantener compacta la trayectoria condensador--MOSFET--inductor--diodo--condensador, con el fin de reducir inductancias parásitas y picos de tensión.

\section{Conclusiones}
\begin{itemize}
\item El diseño ideal del convertidor Boost establece un ciclo de trabajo de aproximadamente \SI{41.2}{\percent} para elevar \SI{10}{V} hasta \SI{17}{V} a \SI{20}{kHz}.
\item Los valores seleccionados de \SI{2.7}{mH} para el inductor y \SI{107}{\micro F} para el condensador permiten alcanzar en simulación una tensión de salida cercana al objetivo.
\item La simulación reportada entrega aproximadamente \SI{14.97}{W} en la carga de \SI{19.3}{ohm} y un rizado de tensión de \SI{0.82}{\percent}, inferior al máximo especificado de \SI{5}{\percent}.
\item El cambio a un MOSFET de canal N y la corrección de la señal PWM fueron necesarios para obtener el comportamiento esperado del convertidor.
\item No se dispone de datos de potencia de entrada, eficiencia, rizado de corriente ni resultados de implementación física; por ello, estas magnitudes quedan pendientes de verificación experimental.
\end{itemize}

\end{document}
